\chapter{Banque de Questions / Réponses et Préparation de l'Épreuve de Spécialité}

Ce chapitre comprend deux volets : (1) une banque de questions/réponses sur le \textbf{ministère et la politique étrangère}, utile pour la rédaction et surtout pour l'oral (coefficient 6) ; (2) la préparation de l'\textbf{épreuve écrite de spécialité} (3 heures, coefficient 4) en informatique et réseaux, avec des questions/réponses par thème et des \textbf{exercices rédigés corrigés} placés dans le contexte du ministère.

\section*{Légende des Niveaux de Priorité et de Typologie}
\begin{itemize}
    \item \textbf{Priorité :}
    \textbf{🔴 Très important} (attribution centrale du ministère ou actualité majeure) ;
    \textbf{🟠 Important} (notion institutionnelle ou diplomatique fondamentale) ;
    \textbf{🟡 À connaître} (notion utile) ;
    \textbf{⚪ Secondaire} (culture générale).
    \item \textbf{Type :}
    \textbf{[Officiel]} = tiré de l'avis de concours ou d'un texte officiel ;
    \textbf{[Actualité]} = fait daté de 2024-2026 ;
    \textbf{[Thème récurrent]} = question classique des jurys de la fonction publique ;
    \textbf{[Prépa]} = question conçue pour ce guide.
\end{itemize}

\section{Volet 1 : Ministère, Diplomatie et Marocains du Monde}

{\footnotesize
\setlength{\tabcolsep}{3pt}
\begin{longtable}{|>{\raggedright\arraybackslash\hspace{0pt}}p{3.9cm}|>{\raggedright\arraybackslash\hspace{0pt}}p{5.1cm}|>{\raggedright\arraybackslash\hspace{0pt}}p{2.1cm}|>{\raggedright\arraybackslash\hspace{0pt}}p{2.3cm}|>{\raggedright\arraybackslash\hspace{0pt}}p{1.5cm}|c|}
\hline
\rowcolor{gray!20} \textbf{Question} & \textbf{Réponse Synthétique} & \textbf{Thème} & \textbf{Source} & \textbf{Type} & \textbf{Prio} \\ \hline
\endfirsthead
\hline
\rowcolor{gray!20} \textbf{Question} & \textbf{Réponse Synthétique} & \textbf{Thème} & \textbf{Source} & \textbf{Type} & \textbf{Prio} \\ \hline
\endhead

Quelles sont les trois épreuves du concours d'ingénieurs d'État du MAECAME 2026 ? & Écrit de spécialité informatique et réseaux (3 h, coef. 4) ; rédaction d'un sujet sur les attributions du ministère (2 h, coef. 2) ; entretien oral (30-40 min, coef. 6). & Concours & Avis C43032/26 & [Officiel] & 🔴 \\ \hline

Combien de postes sont ouverts et dans quelles spécialités ? & 16 postes (échelle 11) : 10 en Réseaux et Télécommunications, 6 en Développement informatique, dont 1 poste dédié au secteur des MRE. & Concours & Avis C43032/26 & [Officiel] & 🔴 \\ \hline

Quelles sont les grandes attributions du ministère ? & Conduire la politique étrangère sous l'autorité du Roi (intégrité territoriale, bilatéral, multilatéral) ; coopération africaine et Sud-Sud ; service et protection des MRE ; diplomatie économique et culturelle ; gestion du réseau diplomatique et consulaire. & Ministère & Décret 2.24.957 & [Officiel] & 🔴 \\ \hline

Qui est le ministre des Affaires étrangères et depuis quand ? & M. Nasser Bourita, depuis le 5 avril 2017, reconduit le 23 octobre 2024 ; gouvernement sortant après les législatives du 23 septembre 2026. & Ministère & Presse officielle & [Actualité] & 🔴 \\ \hline

Quel texte a réorganisé le ministère en 2024 et qu'a-t-il changé ? & Le décret n° 2.24.957 du 21 novembre 2024 : regroupement en pôles, nouvelles directions (dont la Direction des Systèmes d'Information au sein de la DGAAG), transformation de l'AMED en IMFRED. & Ministère & BO / presse & [Officiel] & 🔴 \\ \hline

Quelle est la mission de la Direction des Systèmes d'Information du ministère ? & Élaborer et exécuter les plans du ministère en matière de SI, gérer l'infrastructure et garantir la cybersécurité, pour l'administration centrale et le réseau à l'étranger. & Ministère / SI & Décret 2.24.957 & [Officiel] & 🔴 \\ \hline

Quelle est la taille du réseau diplomatique et consulaire ? & Environ 104 ambassades et missions permanentes et 53 consulats généraux ; 21 nouveaux consuls généraux nommés en 2026 (35 \% des postes). & Ministère & Presse 2026 & [Actualité] & 🟠 \\ \hline

Qu'est-ce que la DACS ? & La Direction des Affaires Consulaires et Sociales : état civil, documents, protection et assistance des Marocains à l'étranger, statistiques consulaires. & Ministère & Organigramme & [Thème récurrent] & 🟠 \\ \hline

Qu'est-ce que l'AMCI ? & L'Agence Marocaine de Coopération Internationale (1986), bras de la coopération Sud-Sud : bourses d'étudiants étrangers (surtout africains), formation, expertise. & Coopération & AMCI & [Thème récurrent] & 🟠 \\ \hline

Qu'est-ce que l'IMFRED ? & L'Institut Marocain de Formation, de Recherche et d'Études Diplomatiques, qui remplace l'Académie Marocaine des Études Diplomatiques (2011) ; forme les diplomates. & Ministère & Décret 2.24.957 & [Officiel] & 🟡 \\ \hline

Quels articles de la Constitution concernent les MRE ? & Articles 16 (protection des droits et intérêts), 17 (droits de citoyenneté, vote), 18 (participation aux institutions consultatives) ; article 163 (CCME). & MRE & Constitution 2011 & [Thème récurrent] & 🔴 \\ \hline

Quelles institutions s'occupent des MRE et que change la réforme de 2024 ? & Ministère (DACS, consulats), CCME (2007), Fondation Hassan II (1990), Fondation Mohammed V (Marhaba). Discours du 6 novembre 2024 : CCME recentré sur la prospective, création d'une Fondation Mohammedia des MRE (bras opérationnel) et d'un mécanisme de mobilisation des compétences ; textes en cours. & MRE & Discours royal 2024 & [Actualité] & 🔴 \\ \hline

Chiffres clés des MRE ? & Plus de 5 millions de personnes ; transferts 117,7 MMDH en 2024 (7,7 \% du PIB), plus de 122 MMDH en 2025, 74,78 MMDH à fin juillet 2026 ; Marhaba 2026 : 4,14 millions accueillis. & MRE & Office des changes / FM5 & [Actualité] & 🔴 \\ \hline

Comment les MRE ont-ils voté aux législatives de 2026 ? & Par procuration, via une plateforme dématérialisée « E-procuration » (24 août -- 22 septembre 2026) ; pas de bureaux de vote à l'étranger. & MRE & FH2MRE & [Actualité] & 🟡 \\ \hline

Que dit la résolution 2797 du Conseil de sécurité ? & Adoptée le 31 octobre 2025 (11 pour, 3 abstentions, l'Algérie ne participant pas) : l'autonomie sous souveraineté marocaine est la base des négociations et « la solution la plus réalisable » ; MINURSO prorogée jusqu'au 31 octobre 2026. & Sahara & ONU & [Actualité] & 🔴 \\ \hline

Quels pays ou organisations soutiennent explicitement le plan d'autonomie ? & États-Unis (10 déc. 2020, réaffirmé avril 2025), Espagne (mars 2022), France (30 juillet 2024), Royaume-Uni (1\textsuperscript{er} juin 2025), Union européenne (position commune du 29 janvier 2026), et une large majorité d'États arabes et africains ; une trentaine de consulats à Laâyoune et Dakhla. & Sahara & Fiche 3 & [Actualité] & 🔴 \\ \hline

Quels sont les six principes énoncés par le Maroc à l'ONU le 22 septembre 2026 ? & Pas de solution hors autonomie sous souveraineté marocaine ; négociation onusienne sur la seule base de l'Initiative ; obligations de la 2797 non négociables ; refus du statu quo ; crédibilité liée au respect du cessez-le-feu ; libération des populations de Tindouf. & Sahara & AGNU 2026 & [Actualité] & 🔴 \\ \hline

Quand le Maroc est-il revenu à l'Union Africaine et pourquoi l'avait-il quittée ? & Retour le 30 janvier 2017 ; retrait de l'OUA en 1984 après l'admission de la « RASD ». Demande d'adhésion à la CEDEAO en 2017, toujours en attente. & Afrique & UA & [Thème récurrent] & 🔴 \\ \hline

Qu'est-ce que l'Initiative Atlantique ? & Initiative royale de novembre 2023 offrant aux pays du Sahel (Mali, Burkina Faso, Niger, Tchad) un accès à l'Atlantique via les infrastructures marocaines, notamment le port Dakhla Atlantique. & Afrique & Discours royal & [Actualité] & 🔴 \\ \hline

Qu'est-ce que le gazoduc Nigeria-Maroc ? & Gazoduc Africain Atlantique d'environ 6 900 km et 25 milliards de dollars traversant 13 pays ; accord intergouvernemental signé par la CEDEAO le 19 juillet 2026 ; premier gaz vers 2031. & Afrique & Presse 2026 & [Actualité] & 🟠 \\ \hline

Quel est le rôle du Maroc sur la question migratoire en Afrique ? & Le Roi est Leader de l'UA sur la migration (2018) ; Pacte de Marrakech (décembre 2018) ; Observatoire Africain des Migrations inauguré à Rabat (18 décembre 2020). & Afrique & UA / ONU & [Thème récurrent] & 🟠 \\ \hline

Qu'est-ce que la Fondation Mohammed VI des Oulémas Africains ? & Fondation créée par dahir en 2015 (siège à Fès) qui fédère les oulémas de 48 sections nationales autour d'un islam du juste milieu ; instrument de la diplomatie religieuse. & Afrique & Habous & [Thème récurrent] & 🟡 \\ \hline

Qu'a décidé la CJUE le 4 octobre 2024 et quelle a été la suite ? & Annulation des accords agricole et de pêche Maroc-UE ; un nouvel accord agricole par échange de lettres a été publié au JOUE le 3 octobre 2025 ; pas de nouvel accord de pêche à ce jour. & UE & CJUE / Consilium & [Actualité] & 🟠 \\ \hline

Quelle est la situation des relations avec l'Algérie ? & Rupture des relations diplomatiques décidée par Alger le 24 août 2021, espace aérien fermé, gazoduc Maghreb-Europe non renouvelé (octobre 2021), frontières fermées depuis 1994 ; politique marocaine de la main tendue. & Voisinage & Discours royaux & [Thème récurrent] & 🟠 \\ \hline

Où en sont les relations Maroc-Israël ? & Déclaration tripartite du 22 décembre 2020 ; le 16 septembre 2026, accord pour élever les bureaux de liaison en ambassades ; le Maroc maintient son soutien à un État palestinien dans les frontières de 1967 et préside le Comité Al-Qods. & Moyen-Orient & Presse 2026 & [Actualité] & 🟠 \\ \hline

Quels engagements récents du Maroc pour la paix ? & Accord de participation à la Force internationale de stabilisation à Gaza (Rabat, 15 juillet 2026) ; membre fondateur du Conseil de paix ; contributeur historique aux opérations de maintien de la paix ; accord de Skhirat (Libye, 2015). & Multilatéral & Presse 2026 & [Actualité] & 🟡 \\ \hline

Le Maroc et le Conseil des droits de l'homme ? & Membre 2023-2025 ; présidence en 2024 (ambassadeur Omar Zniber). & Multilatéral & ONU & [Actualité] & 🟡 \\ \hline

Le Maroc est-il candidat au Conseil de sécurité ? & Oui, pour un siège non permanent en 2028-2029 (dernier mandat 2012-2013), candidature soutenue par le Sommet arabe de Bagdad (mai 2025) et confirmée à l'AGNU le 22 septembre 2026. & Multilatéral & AGNU 2026 & [Actualité] & 🟠 \\ \hline

Quand et à qui a été attribuée la Coupe du monde 2030 ? & Le 11 décembre 2024 au Maroc, à l'Espagne et au Portugal (13 juin -- 21 juillet 2030) ; levier de diplomatie économique et sportive. & Diplomatie économique & FIFA & [Actualité] & 🟠 \\ \hline

Qu'est-ce que la diplomatie économique et qui la porte au ministère ? & L'action extérieure de l'État pour attirer les IDE, promouvoir les exportations et accompagner les entreprises ; Direction de la Diplomatie Économique, ambassades, en lien avec l'AMDIE et les CRI ; l'un des quatre programmes budgétaires du ministère. & Diplomatie économique & Budget 2026 & [Thème récurrent] & 🟠 \\ \hline

Quels sont les quatre programmes du budget du ministère ? & Modernisation de l'outil diplomatique ; diplomatie économique ; action bilatérale et multilatérale ; services consulaires et état civil. 155 postes budgétaires créés en 2026. & Ministère & Présentation budget & [Actualité] & 🟡 \\ \hline

Qu'est-ce que l'e-Visa marocain ? & Visa électronique lancé le 10 juillet 2022 sur acces-maroc.ma (72 h, ou 24 h en express, valable 30 jours) ; 385 738 délivrés à fin 2024, 111 nationalités. & Numérique & diplomatie.ma & [Actualité] & 🔴 \\ \hline

Quels services consulaires sont numérisés ? & consulat.ma (guide), rdv.consulat.ma (rendez-vous), eConsulat (plateforme interne, nouveaux modules), centre d'appel multilingue, e-timbre et e-paiement, apostille.ma, e-procuration ; bureau d'ordre digital et parapheur électronique en interne. & Numérique & Presse 2025 & [Actualité] & 🔴 \\ \hline

Quel cadre juridique encadre la cybersécurité et les données au Maroc ? & Loi 05-20 (cybersécurité) et décret 2-21-406 ; DGSSI et DNSSI ; Stratégie nationale de cybersécurité 2030 (juillet 2024) ; loi 09-08 et CNDP (données personnelles) ; loi 43-20 (services de confiance). & Numérique & DGSSI / CNDP & [Thème récurrent] & 🔴 \\ \hline

Que prévoient les Conventions de Vienne de 1961 et 1963 ? & 1961 : relations diplomatiques (ambassades, immunités, inviolabilité de la valise diplomatique) ; 1963 : relations consulaires (fonctions et protection consulaires). & Droit international & Conventions de Vienne & [Thème récurrent] & 🟡 \\ \hline

Quelle est la différence entre une ambassade et un consulat ? & L'ambassade représente l'État auprès d'un autre État (relations politiques) ; le consulat sert les ressortissants (état civil, documents, protection) et délivre les visas dans sa circonscription. & Droit international & Conventions de Vienne & [Thème récurrent] & 🟡 \\ \hline

Qu'est-ce que l'apostille ? & Certification simplifiée des documents publics entre États parties à la Convention de La Haye de 1961 ; en vigueur au Maroc depuis le 14 août 2016 ; demande en ligne sur apostille.ma. & Consulaire & HCCH & [Thème récurrent] & ⚪ \\ \hline

Quelles sont les dates des grands discours royaux à connaître ? & Discours du Trône (30 juillet), de la Révolution du Roi et du Peuple (20 août), de la Marche Verte (6 novembre) ; discours d'Addis-Abeba (31 janvier 2017) ; discours au Parlement (octobre). & Culture & maroc.ma & [Thème récurrent] & ⚪ \\ \hline
\end{longtable}
}

\section{Volet 2 : Épreuve de Spécialité --- Informatique et Réseaux (3 h, coefficient 4)}

\begin{tcolorbox}[colback=red!4!white,colframe=red!60!black,title=\textbf{Comment se préparer sans annale}]
L'avis précise seulement que l'épreuve « concerne la ou les spécialités demandées, les fonctions ou les postes à occuper en informatique et en réseaux ». Le format peut être un QCM, des questions ouvertes ou des exercices. Les blocs A à I ci-dessous couvrent le socle (format question $\rightarrow$ réponse, sigles expliqués entre parenthèses) ; le bloc J ajoute les \textbf{réseaux et télécommunications approfondis} (10 postes sur 16) ; le bloc K propose des \textbf{exercices rédigés corrigés} dans le contexte du ministère, à refaire sur papier.
\end{tcolorbox}

\subsection*{A. Algorithmique et Programmation}
\begin{enumerate}[leftmargin=*]
    \item \textbf{Complexité d'une recherche dichotomique (on coupe le tableau trié en deux à chaque étape) sur $n$ éléments ?} $\rightarrow$ $O(\log n)$ (la notation « grand O » mesure la croissance du temps de calcul quand $n$ augmente). Piège : $O(n)$ est la recherche séquentielle (élément par élément).
    \item \textbf{Complexité moyenne du tri rapide (\textit{quicksort}) et du tri fusion ?} $\rightarrow$ $O(n \log n)$ ; le tri rapide peut dégénérer en $O(n^2)$ dans le pire cas, le tri fusion reste en $O(n \log n)$. Le tri à bulles est en $O(n^2)$.
    \item \textbf{Pile ou file : laquelle suit le principe LIFO ?} $\rightarrow$ La \textbf{pile} (LIFO = \textit{Last In, First Out} : dernier entré, premier sorti). La \textbf{file} suit le principe FIFO (\textit{First In, First Out} : premier entré, premier sorti).
    \item \textbf{Quelle structure de données donne un accès moyen en $O(1)$ (temps constant) par clé ?} $\rightarrow$ La \textbf{table de hachage} (dictionnaire).
    \item \textbf{Quelle est la hauteur d'un arbre binaire de recherche équilibré contenant $n$ nœuds ?} $\rightarrow$ Environ $\log_2 n$.
    \item \textbf{Parcours en largeur (BFS) ou en profondeur (DFS) : lequel utilise une file ?} $\rightarrow$ Le parcours en \textbf{largeur} (BFS = \textit{Breadth-First Search}) ; le parcours en profondeur (DFS = \textit{Depth-First Search}) utilise une pile ou la récursivité.
    \item \textbf{Quel algorithme calcule le plus court chemin dans un graphe à poids positifs ?} $\rightarrow$ L'algorithme de \textbf{Dijkstra}.
    \item \textbf{Qu'est-ce qu'une fonction récursive ?} $\rightarrow$ Une fonction qui s'appelle elle-même ; elle doit avoir une \textbf{condition d'arrêt}, sinon débordement de pile (\textit{stack overflow}).
    \item \textbf{Combien vaut 1011 en binaire, en décimal ?} $\rightarrow$ $8 + 0 + 2 + 1 = $ \textbf{11} ; en hexadécimal (base 16), 11 s'écrit B.
    \item \textbf{Langage compilé ou interprété ?} $\rightarrow$ C et C++ sont compilés en code machine ; Python est interprété ; Java est compilé en \textit{bytecode} exécuté par la JVM (\textit{Java Virtual Machine}, machine virtuelle Java).
\end{enumerate}

\subsection*{B. Programmation Orientée Objet (POO) et Génie Logiciel}
\begin{enumerate}[leftmargin=*]
    \item \textbf{Quels sont les principes fondamentaux de la POO ?} $\rightarrow$ \textbf{Encapsulation} (cacher les données derrière des méthodes), \textbf{héritage} (une classe fille réutilise une classe mère), \textbf{polymorphisme} (une même méthode se comporte différemment selon l'objet), \textbf{abstraction}.
    \item \textbf{Surcharge ou redéfinition ?} $\rightarrow$ La \textbf{surcharge} (\textit{overloading}) = même nom de méthode, paramètres différents, dans la même classe ; la \textbf{redéfinition} (\textit{overriding}) = la classe fille réécrit une méthode de la classe mère.
    \item \textbf{Interface ou classe abstraite en Java ?} $\rightarrow$ Une classe ne peut hériter que d'\textbf{une seule} classe (abstraite ou non) mais peut implémenter \textbf{plusieurs} interfaces.
    \item \textbf{Que signifie le mot-clé \texttt{static} ?} $\rightarrow$ L'attribut ou la méthode appartient à la classe, et non à chaque objet.
    \item \textbf{Que signifie SOLID ?} $\rightarrow$ 5 principes de conception : responsabilité unique (\textit{Single responsibility}), ouvert/fermé (\textit{Open/closed}), substitution de Liskov, ségrégation des interfaces, inversion des dépendances.
    \item \textbf{Quel patron de conception (\textit{design pattern}) garantit une seule instance d'une classe ?} $\rightarrow$ Le \textbf{Singleton}. Autres classiques : Factory (fabrique d'objets), Observer (notification des abonnés), MVC (Modèle-Vue-Contrôleur : séparer données, affichage et logique).
    \item \textbf{Cycle en V ou méthode agile ?} $\rightarrow$ Le \textbf{cycle en V} enchaîne des phases planifiées (spécification $\rightarrow$ conception $\rightarrow$ codage $\rightarrow$ tests de chaque niveau) ; l'\textbf{agile} livre par petites itérations avec le client.
    \item \textbf{Quelle méthode agile s'organise en sprints avec un Product Owner et un Scrum Master ?} $\rightarrow$ \textbf{Scrum} (sprint = itération de 1 à 4 semaines ; Product Owner = représentant du métier qui priorise le \textit{backlog}, la liste des besoins ; Scrum Master = facilitateur de l'équipe).
    \item \textbf{Tests unitaires, d'intégration, de recette ?} $\rightarrow$ Unitaires : une fonction isolée ; intégration : plusieurs modules ensemble ; recette (UAT, \textit{User Acceptance Testing}) : validation par l'utilisateur final.
    \item \textbf{Que signifie CI/CD ?} $\rightarrow$ \textit{Continuous Integration / Continuous Delivery-Deployment} (intégration continue / livraison ou déploiement continu) : automatiser la compilation, les tests et la mise en production (ex. GitLab CI, Jenkins).
    \item \textbf{À quoi sert Git ?} $\rightarrow$ Gestion de versions du code : \texttt{commit} (enregistrer), \texttt{branch} (branche de travail parallèle), \texttt{merge} (fusionner), \texttt{pull request} (demande de relecture avant fusion).
    \item \textbf{Qu'est-ce qu'une API REST ?} $\rightarrow$ API = \textit{Application Programming Interface} (interface pour que deux logiciels communiquent) ; REST = \textit{Representational State Transfer} : interface web \textbf{sans état} (le serveur ne garde pas la session) qui manipule des ressources identifiées par des URL avec les méthodes HTTP : GET (lire), POST (créer), PUT (modifier), DELETE (supprimer). Données souvent en JSON.
    \item \textbf{Codes HTTP à connaître ?} $\rightarrow$ 200 (succès), 201 (créé), 301 (redirection), 401 (non authentifié), 403 (interdit), 404 (introuvable), 500 (erreur serveur).
    \item \textbf{Architecture monolithique ou microservices ?} $\rightarrow$ Monolithe : une seule application ; microservices : petits services indépendants communiquant par API, déployables séparément (plus souples, mais plus complexes à superviser).
    \item \textbf{Quel diagramme UML (\textit{Unified Modeling Language}, langage de modélisation) représente les interactions entre objets dans l'ordre chronologique ?} $\rightarrow$ Le \textbf{diagramme de séquence}. Autres : cas d'utilisation (qui fait quoi avec le système), classes (structure), activité (déroulement d'un processus).
\end{enumerate}

\subsection*{C. Bases de Données et SQL}
\begin{enumerate}[leftmargin=*]
    \item \textbf{Que signifie ACID ?} $\rightarrow$ Les 4 garanties d'une transaction : \textbf{Atomicité} (tout ou rien), \textbf{Cohérence} (la base reste valide), \textbf{Isolation} (les transactions simultanées ne se gênent pas), \textbf{Durabilité} (une fois validée, la donnée n'est pas perdue).
    \item \textbf{Clé primaire et clé étrangère ?} $\rightarrow$ La clé primaire identifie de façon unique chaque ligne ; la clé étrangère fait référence à la clé primaire d'une autre table (lien entre tables).
    \item \textbf{Quelle forme normale élimine les dépendances partielles ?} $\rightarrow$ La \textbf{2FN} (FN = Forme Normale : règle qui évite les doublons et incohérences). 1FN : valeurs atomiques (une seule valeur par case) ; 2FN : 1FN + chaque attribut dépend de \textbf{toute} la clé ; 3FN : 2FN + aucun attribut ne dépend d'un autre attribut non-clé (pas de dépendance transitive).
    \item \textbf{Quelle méthode française de conception utilise le MCD et le MLD ?} $\rightarrow$ \textbf{Merise} (MCD = Modèle Conceptuel de Données, entités et associations ; MLD = Modèle Logique de Données, tables ; MPD = Modèle Physique de Données).
    \item \textbf{Une association « plusieurs à plusieurs » (n:n) entre deux tables se traduit par :} $\rightarrow$ une \textbf{table de jointure} (table intermédiaire contenant les deux clés étrangères).
    \item \textbf{Quelle requête SQL regroupe des lignes pour appliquer une fonction d'agrégat (SUM, COUNT, AVG) ?} $\rightarrow$ \texttt{GROUP BY}. Piège : on filtre les groupes avec \texttt{HAVING}, pas avec \texttt{WHERE} (qui filtre les lignes avant regroupement).
    \item \textbf{\texttt{INNER JOIN} ou \texttt{LEFT JOIN} ?} $\rightarrow$ INNER : seulement les lignes qui correspondent dans les deux tables ; LEFT : toutes les lignes de la table de gauche, même sans correspondance (valeurs NULL).
    \item \textbf{Catégories du langage SQL ?} $\rightarrow$ LDD/DDL (définition : CREATE, ALTER, DROP), LMD/DML (manipulation : SELECT, INSERT, UPDATE, DELETE), LCD/DCL (droits : GRANT, REVOKE), TCL (transactions : COMMIT, ROLLBACK).
    \item \textbf{\texttt{DELETE}, \texttt{TRUNCATE} ou \texttt{DROP} ?} $\rightarrow$ DELETE supprime des lignes (avec condition possible) ; TRUNCATE vide toute la table ; DROP supprime la table elle-même.
    \item \textbf{À quoi sert un index ?} $\rightarrow$ Accélérer les recherches (comme l'index d'un livre), au prix d'écritures un peu plus lentes.
    \item \textbf{Qu'est-ce qu'une vue ?} $\rightarrow$ Une table virtuelle définie par une requête SELECT (utile pour simplifier et restreindre l'accès aux données).
    \item \textbf{SQL ou NoSQL ?} $\rightarrow$ SQL : tables relationnelles, schéma fixe, transactions ACID (Oracle, PostgreSQL, MySQL) ; NoSQL (\textit{Not only SQL}) : documents, clé-valeur, graphes, schéma souple et forte montée en charge (MongoDB, Redis, Cassandra).
    \item \textbf{Que dit le théorème CAP ?} $\rightarrow$ Un système distribué ne peut garantir à la fois que \textbf{deux} des trois propriétés : Cohérence (\textit{Consistency}), Disponibilité (\textit{Availability}), tolérance au partitionnement (\textit{Partition tolerance}).
\end{enumerate}

\subsection*{D. Systèmes d'Exploitation, Virtualisation et Cloud}
\begin{enumerate}[leftmargin=*]
    \item \textbf{Processus ou thread ?} $\rightarrow$ Un processus est un programme en cours d'exécution avec sa propre mémoire ; un thread (fil d'exécution) partage la mémoire du processus.
    \item \textbf{Qu'est-ce qu'un interblocage (\textit{deadlock}) ?} $\rightarrow$ Deux processus s'attendent mutuellement pour une ressource et restent bloqués indéfiniment.
    \item \textbf{Commandes Linux utiles ?} $\rightarrow$ \texttt{ls} (lister), \texttt{cd} (changer de dossier), \texttt{chmod} (droits), \texttt{chown} (propriétaire), \texttt{ps}/\texttt{top} (processus), \texttt{grep} (rechercher du texte), \texttt{sudo} (exécuter en administrateur).
    \item \textbf{Que signifie \texttt{chmod 755} ?} $\rightarrow$ Propriétaire : lecture + écriture + exécution (7 = 4+2+1) ; groupe et autres : lecture + exécution (5 = 4+1).
    \item \textbf{Mémoire virtuelle ?} $\rightarrow$ Technique qui utilise le disque pour étendre la mémoire vive (RAM, \textit{Random Access Memory}) grâce à la pagination.
    \item \textbf{Conteneur (Docker) ou machine virtuelle ?} $\rightarrow$ Le conteneur partage le noyau (cœur) du système hôte, il est donc plus léger et démarre vite ; la VM (\textit{Virtual Machine}) embarque son propre système d'exploitation via un \textbf{hyperviseur} (logiciel qui fait tourner plusieurs VM : VMware, Hyper-V, KVM).
    \item \textbf{Rôle de Kubernetes ?} $\rightarrow$ Orchestrer (déployer, répartir, redémarrer, mettre à l'échelle) des conteneurs sur plusieurs serveurs.
    \item \textbf{Les trois modèles de service du cloud ?} $\rightarrow$ \textbf{IaaS} (\textit{Infrastructure as a Service} : on loue serveurs et stockage), \textbf{PaaS} (\textit{Platform as a Service} : on loue une plateforme pour déployer son code), \textbf{SaaS} (\textit{Software as a Service} : on utilise directement un logiciel en ligne, ex. messagerie).
    \item \textbf{Modèles de déploiement du cloud ?} $\rightarrow$ Public, privé, hybride (mélange des deux), communautaire. Pour les données sensibles de l'État : cloud privé ou \textbf{souverain} (hébergé au Maroc).
    \item \textbf{Scalabilité verticale ou horizontale ?} $\rightarrow$ Verticale : ajouter de la puissance à une machine ; horizontale : ajouter des machines.
\end{enumerate}

\subsection*{E. Réseaux et Systèmes}
\begin{enumerate}[leftmargin=*]
    \item \textbf{Combien de couches compte le modèle OSI (\textit{Open Systems Interconnection}, modèle de référence des réseaux) ?} $\rightarrow$ \textbf{7} : Physique, Liaison, Réseau, Transport, Session, Présentation, Application. Le modèle TCP/IP en compte 4.
    \item \textbf{À quelle couche travaille un routeur ? un commutateur (switch) ?} $\rightarrow$ Routeur : couche 3 (Réseau, adresses IP = adresses logiques) ; switch : couche 2 (Liaison, adresses MAC = adresse physique unique de la carte réseau).
    \item \textbf{TCP ou UDP ?} $\rightarrow$ TCP (\textit{Transmission Control Protocol}) : avec connexion, fiable (accusés de réception) ; UDP (\textit{User Datagram Protocol}) : sans connexion, plus rapide, sans garantie (vidéo, voix, DNS).
    \item \textbf{Combien d'adresses hôtes utilisables dans un réseau /24 ?} $\rightarrow$ /24 signifie que 24 bits sur 32 désignent le réseau (masque 255.255.255.0) ; il reste 8 bits : $2^8 - 2 = $ \textbf{254} (on retire l'adresse du réseau et l'adresse de diffusion). Formule : $2^{(32 - n)} - 2$.
    \item \textbf{Quelle plage d'adresses IPv4 est privée ?} $\rightarrow$ 10.0.0.0/8, 172.16.0.0/12, 192.168.0.0/16 (non routables sur Internet, traduites par le \textbf{NAT}, \textit{Network Address Translation}).
    \item \textbf{Quel protocole attribue automatiquement les adresses IP ?} $\rightarrow$ \textbf{DHCP} (\textit{Dynamic Host Configuration Protocol}) ; le \textbf{DNS} (\textit{Domain Name System}) traduit un nom (finances.gov.ma) en adresse IP ; \textbf{ARP} trouve l'adresse MAC correspondant à une IP.
    \item \textbf{Ports par défaut ?} $\rightarrow$ HTTP 80, HTTPS 443, SSH 22 (connexion sécurisée à distance), FTP 21 (transfert de fichiers), DNS 53, SMTP 25 (envoi de courriel), RDP 3389 (bureau à distance Windows).
    \item \textbf{Quelle commande teste la joignabilité d'une machine ?} $\rightarrow$ \texttt{ping} (protocole ICMP) ; \texttt{traceroute}/\texttt{tracert} affiche le chemin suivi.
    \item \textbf{Taille d'une adresse IPv6 ?} $\rightarrow$ \textbf{128 bits} (IPv4 : 32 bits).
    \item \textbf{À quoi sert un VLAN (\textit{Virtual LAN}, réseau local virtuel) ?} $\rightarrow$ Séparer logiquement un même réseau physique en plusieurs réseaux isolés (ex. un VLAN par direction).
    \item \textbf{Qu'est-ce qu'un VPN (\textit{Virtual Private Network}) ?} $\rightarrow$ Un tunnel chiffré qui relie de façon sécurisée un utilisateur ou un site distant au réseau de l'organisation (ex. IPsec, SSL/TLS).
    \item \textbf{Qu'est-ce qu'une DMZ (zone démilitarisée) ?} $\rightarrow$ Une zone réseau intermédiaire, isolée par pare-feu, où l'on place les serveurs accessibles depuis Internet (site web, messagerie).
    \item \textbf{Quel niveau RAID (\textit{Redundant Array of Independent Disks}, regroupement de disques) utilise une parité répartie ?} $\rightarrow$ \textbf{RAID 5} (supporte la panne d'un disque). RAID 0 = vitesse sans redondance ; RAID 1 = miroir (copie identique).
    \item \textbf{Que signifient PRA et PCA ?} $\rightarrow$ Plan de Reprise d'Activité (redémarrer après un sinistre) et Plan de Continuité d'Activité (continuer à fonctionner pendant l'incident). Indicateurs : RTO (durée maximale d'interruption acceptable) et RPO (perte de données maximale acceptable).
    \item \textbf{Qu'est-ce que l'annuaire Active Directory / LDAP ?} $\rightarrow$ Un annuaire centralisé des utilisateurs et des droits (LDAP = \textit{Lightweight Directory Access Protocol}).
\end{enumerate}

\subsection*{F. Cybersécurité}
\begin{enumerate}[leftmargin=*]
    \item \textbf{Les trois piliers de la sécurité (triade DIC, en anglais CIA) ?} $\rightarrow$ \textbf{D}isponibilité (le service est accessible), \textbf{I}ntégrité (la donnée n'est pas modifiée), \textbf{C}onfidentialité (seules les personnes autorisées y accèdent). On ajoute souvent la \textbf{traçabilité} (preuve de qui a fait quoi).
    \item \textbf{Quel texte encadre la cybersécurité au Maroc ?} $\rightarrow$ La \textbf{loi n° 05-20} (2020), sous l'autorité nationale \textbf{DGSSI} ; la DNSSI (Directive Nationale de la Sécurité des SI) fixe les règles minimales des administrations.
    \item \textbf{Autres textes du numérique marocain ?} $\rightarrow$ Loi 09-08 (données personnelles, CNDP), loi 43-20 (services de confiance : signature et cachet électroniques), loi 07-03 (atteintes aux systèmes de traitement automatisé des données, dans le Code pénal), loi 55-19 (simplification des procédures administratives).
    \item \textbf{Quelle norme définit un SMSI (Système de Management de la Sécurité de l'Information) ?} $\rightarrow$ \textbf{ISO/IEC 27001} ; ISO 27005 porte sur la gestion des risques (méthode française équivalente : EBIOS).
    \item \textbf{Chiffrement symétrique ou asymétrique ?} $\rightarrow$ \textbf{Symétrique} : la même clé chiffre et déchiffre (AES = \textit{Advanced Encryption Standard}, rapide). \textbf{Asymétrique} : paire clé publique (pour chiffrer) / clé privée (pour déchiffrer) (RSA = initiales de ses inventeurs Rivest, Shamir et Adleman).
    \item \textbf{Signature électronique : avec quelle clé signe-t-on ?} $\rightarrow$ Avec sa \textbf{clé privée} ; le destinataire vérifie avec la \textbf{clé publique}. Elle garantit l'authenticité, l'intégrité et la non-répudiation (l'auteur ne peut pas nier). Les clés publiques sont certifiées par une PKI (\textit{Public Key Infrastructure}, autorité de certification).
    \item \textbf{Une fonction de hachage (SHA-256) garantit :} $\rightarrow$ l'\textbf{intégrité} (une empreinte unique de la donnée). Piège : elle ne garantit pas la confidentialité, car le hachage n'est pas réversible (on ne « déchiffre » pas un hash).
    \item \textbf{Comment stocker les mots de passe ?} $\rightarrow$ Hachés avec un \textbf{sel} (valeur aléatoire) et un algorithme lent (bcrypt, Argon2), jamais en clair.
    \item \textbf{Qu'est-ce que la MFA ?} $\rightarrow$ \textit{Multi-Factor Authentication} (authentification multifacteur) : combiner au moins deux éléments parmi ce que l'on sait (mot de passe), ce que l'on possède (téléphone, carte) et ce que l'on est (empreinte).
    \item \textbf{Hameçonnage (phishing) ?} $\rightarrow$ Faux message imitant un organisme de confiance pour voler des identifiants ; parades : sensibilisation, MFA, filtrage des courriels.
    \item \textbf{Rançongiciel (ransomware) ?} $\rightarrow$ Logiciel qui chiffre les données et exige une rançon ; parades : sauvegardes hors ligne (règle 3-2-1 : 3 copies, 2 supports, 1 hors site), mises à jour, segmentation du réseau.
    \item \textbf{DDoS ?} $\rightarrow$ \textit{Distributed Denial of Service} (déni de service distribué) : saturer un service depuis de nombreuses machines pour le rendre indisponible.
    \item \textbf{Injection SQL et XSS ?} $\rightarrow$ Injection SQL : insérer du code SQL via un champ de saisie (parade : requêtes paramétrées) ; XSS (\textit{Cross-Site Scripting}) : injecter du JavaScript dans une page vue par d'autres (parade : échapper les sorties). Référence : OWASP Top 10 (liste des 10 risques web majeurs).
    \item \textbf{Attaque de l'homme du milieu (MITM) ?} $\rightarrow$ L'attaquant s'intercale entre deux parties pour lire ou modifier les échanges ; parade : chiffrement TLS (le « S » de HTTPS) et certificats.
    \item \textbf{Principe du « moindre privilège » et Zero Trust ?} $\rightarrow$ Donner uniquement les droits nécessaires ; Zero Trust : ne faire confiance à aucun accès par défaut, tout vérifier à chaque fois.
    \item \textbf{Pare-feu, IDS, IPS ?} $\rightarrow$ Pare-feu (\textit{firewall}) : filtre le trafic selon des règles ; IDS (\textit{Intrusion Detection System}) : détecte et alerte ; IPS (\textit{Intrusion Prevention System}) : détecte et bloque.
    \item \textbf{Qu'est-ce qu'un SOC et un SIEM ?} $\rightarrow$ SOC (\textit{Security Operations Center}) : équipe qui surveille la sécurité en continu ; SIEM (\textit{Security Information and Event Management}) : outil qui centralise et analyse les journaux (\textit{logs}).
    \item \textbf{Qui protège les données personnelles au Maroc ?} $\rightarrow$ La \textbf{CNDP} (loi 09-08). Piège : la DGSSI gère la sécurité des systèmes, pas les droits des personnes.
    \item \textbf{Qu'est-ce qu'un test d'intrusion (pentest) ?} $\rightarrow$ Une attaque simulée et autorisée pour découvrir les failles avant les pirates.
\end{enumerate}

\subsection*{G. Systèmes Décisionnels, Data Science et IA}
\begin{enumerate}[leftmargin=*]
    \item \textbf{Que signifie ETL ?} $\rightarrow$ \textit{Extract, Transform, Load} (extraire les données des sources, les nettoyer et les transformer, les charger dans l'entrepôt de données ou \textit{data warehouse}).
    \item \textbf{OLTP ou OLAP ?} $\rightarrow$ OLTP (\textit{Online Transaction Processing}) : transactions du quotidien (ex. saisie d'une dépense dans GID) ; OLAP (\textit{Online Analytical Processing}) : analyse multidimensionnelle pour décider (cubes, tableaux de bord).
    \item \textbf{Schéma en étoile ?} $\rightarrow$ Une table de \textbf{faits} (ce qu'on mesure : montants, quantités) au centre, reliée à des tables de \textbf{dimensions} (axes d'analyse : temps, région, contribuable). Variante normalisée : schéma en flocon.
    \item \textbf{Data warehouse ou data lake ?} $\rightarrow$ Entrepôt : données structurées et nettoyées pour l'analyse ; lac de données : données brutes de tous formats stockées telles quelles.
    \item \textbf{Qu'est-ce qu'un KPI ?} $\rightarrow$ \textit{Key Performance Indicator} (indicateur clé de performance), ex. délai moyen de paiement des fournisseurs de l'État. Outils de visualisation : Power BI, Tableau.
    \item \textbf{Apprentissage supervisé ou non supervisé ?} $\rightarrow$ Supervisé : données étiquetées (on connaît la réponse), pour la classification (catégorie) ou la régression (valeur numérique) ; non supervisé : sans étiquette, ex. clustering (regroupement automatique, k-means).
    \item \textbf{Quel indicateur choisir pour détecter une fraude rare ?} $\rightarrow$ Le \textbf{rappel} (part des fraudes réellement détectées), la \textbf{précision} (part des alertes qui sont de vraies fraudes) ou le F1-score (moyenne des deux). Piège : l'exactitude (\textit{accuracy}) est trompeuse quand les fraudes sont très rares.
    \item \textbf{Qu'est-ce qu'une matrice de confusion ?} $\rightarrow$ Un tableau qui compare prédictions et réalité : vrais positifs, faux positifs, vrais négatifs, faux négatifs.
    \item \textbf{Surapprentissage (overfitting) ?} $\rightarrow$ Le modèle apprend « par cœur » les données d'entraînement et se trompe sur de nouvelles données ; parades : validation croisée (tester sur des données mises de côté), régularisation, plus de données.
    \item \textbf{Pourquoi séparer jeu d'entraînement et jeu de test ?} $\rightarrow$ Pour évaluer le modèle sur des données qu'il n'a jamais vues (souvent 80 \% / 20 \%).
    \item \textbf{Que sont les « 3V » du Big Data ?} $\rightarrow$ \textbf{Volume} (quantité), \textbf{Vélocité} (vitesse de production), \textbf{Variété} (formats différents), souvent complétés par Véracité (fiabilité) et Valeur. Outils : Hadoop, Spark.
    \item \textbf{Qu'est-ce qu'un réseau de neurones profond (\textit{deep learning}) et un LLM ?} $\rightarrow$ Un modèle composé de nombreuses couches de neurones artificiels ; un LLM (\textit{Large Language Model}, grand modèle de langage) est un modèle de ce type entraîné sur d'immenses textes pour comprendre et générer du langage (IA générative).
    \item \textbf{Risques de l'IA dans l'administration ?} $\rightarrow$ Biais (discriminations), manque d'explicabilité, protection des données (loi 09-08), dépendance technologique : d'où la nécessité d'une gouvernance et d'un contrôle humain.
    \item \textbf{Exemples d'usage dans l'administration ?} $\rightarrow$ Scoring (notation du risque) des déclarations fiscales pour cibler les contrôles (DGI), ciblage des conteneurs à risque (ADII), détection de fraude documentaire et priorisation des demandes de visa (MAECAME), tableaux de bord des services consulaires.
\end{enumerate}

\subsection*{H. Gouvernance des SI et Gestion de Projet}
\begin{enumerate}[leftmargin=*]
    \item \textbf{ITIL ?} $\rightarrow$ \textit{Information Technology Infrastructure Library} : bonnes pratiques de gestion des services informatiques (gestion des incidents, des problèmes, des changements, catalogue de services).
    \item \textbf{COBIT ?} $\rightarrow$ Référentiel de \textbf{gouvernance} et d'audit des SI (aligner l'informatique sur les objectifs de l'organisation et maîtriser les risques).
    \item \textbf{TOGAF et urbanisation du SI ?} $\rightarrow$ TOGAF (\textit{The Open Group Architecture Framework}) : méthode d'architecture d'entreprise ; l'\textbf{urbanisation} organise le SI comme une ville en zones et quartiers pour le faire évoluer sans tout reconstruire.
    \item \textbf{Qu'est-ce qu'un SLA ?} $\rightarrow$ \textit{Service Level Agreement} (accord de niveau de service) : engagement chiffré d'un prestataire (disponibilité 99,9 \%, délai de résolution).
    \item \textbf{MOA ou MOE ?} $\rightarrow$ Maîtrise d'Ouvrage (MOA) : le métier qui exprime le besoin et valide ; Maîtrise d'Œuvre (MOE) : l'équipe technique qui réalise.
    \item \textbf{Qu'est-ce qu'un cahier des charges et un CPS ?} $\rightarrow$ Le cahier des charges décrit le besoin ; dans un marché public, le CPS (Cahier des Prescriptions Spéciales) fixe les clauses propres au marché (décret 2-22-431 sur les marchés publics).
    \item \textbf{Diagramme de Gantt et chemin critique (PERT) ?} $\rightarrow$ Gantt : planning des tâches sur une ligne de temps ; le chemin critique est la suite de tâches sans marge : tout retard sur elles retarde le projet.
    \item \textbf{Triangle d'un projet ?} $\rightarrow$ \textbf{Coût, Délai, Qualité (périmètre)} : modifier l'un impacte les autres.
    \item \textbf{Qu'est-ce que le TCO ?} $\rightarrow$ \textit{Total Cost of Ownership} (coût total de possession) : achat + déploiement + maintenance + formation + exploitation sur toute la durée de vie. Critère clé pour l'argent public.
    \item \textbf{ERP ?} $\rightarrow$ \textit{Enterprise Resource Planning} (progiciel de gestion intégré) : un seul logiciel pour la finance, les achats, les RH\ldots{} Un ERP administratif centralise par exemple les ressources humaines et le budget d'un ministère.
    \item \textbf{Interopérabilité ?} $\rightarrow$ Capacité de systèmes différents à échanger des données via des standards et des API (ex. échanges DGI–CNSS–RSU), condition du principe « dites-le nous une fois ».
    \item \textbf{Open data ?} $\rightarrow$ Publication de données publiques réutilisables librement (portail data.gov.ma).
\end{enumerate}

\subsection*{I. Compléments Rapides (Pièges Fréquents des QCM de Concours)}
\begin{enumerate}[leftmargin=*]
    \item \textbf{Que signifie CRUD ?} $\rightarrow$ \textit{Create, Read, Update, Delete} (créer, lire, modifier, supprimer) : les 4 opérations de base sur les données (en SQL : INSERT, SELECT, UPDATE, DELETE).
    \item \textbf{Système de fichiers par défaut sous Linux ? sous Windows ?} $\rightarrow$ \textbf{ext4} (Linux) ; \textbf{NTFS} (Windows).
    \item \textbf{Quel protocole sécurise les réseaux Wi-Fi modernes ?} $\rightarrow$ \textbf{WPA2}, et désormais \textbf{WPA3} (WPA = \textit{Wi-Fi Protected Access}). Piège : WEP est obsolète et cassable.
    \item \textbf{Routage statique ou dynamique ?} $\rightarrow$ Statique : routes saisies à la main ; dynamique : les routeurs s'échangent les routes via un protocole (OSPF à l'intérieur d'une organisation, BGP entre opérateurs sur Internet).
    \item \textbf{Quelle commande interroge un serveur DNS ?} $\rightarrow$ \texttt{nslookup} (ou \texttt{dig} sous Linux) ; \texttt{ipconfig} (Windows) / \texttt{ip a} (Linux) affichent la configuration IP.
    \item \textbf{Autres ports à connaître ?} $\rightarrow$ POP3 110 et IMAP 143 (réception de courriels), MySQL 3306, SFTP 22 (transfert de fichiers via SSH).
    \item \textbf{Qu'est-ce qu'une attaque CSRF ?} $\rightarrow$ \textit{Cross-Site Request Forgery} : faire exécuter à un utilisateur connecté une action qu'il n'a pas voulue ; parade : jeton anti-CSRF.
    \item \textbf{Qu'est-ce qu'un certificat X.509 ?} $\rightarrow$ Le format standard des certificats numériques qui lient une clé publique à une identité (utilisés par HTTPS/TLS, TLS = \textit{Transport Layer Security}).
    \item \textbf{En Java, \texttt{==} ou \texttt{equals()} pour comparer deux chaînes ?} $\rightarrow$ \texttt{equals()} compare le contenu ; \texttt{==} compare les références (l'adresse en mémoire). Piège classique.
    \item \textbf{\texttt{git merge} ou \texttt{git rebase} ?} $\rightarrow$ merge fusionne en créant un commit de fusion ; rebase rejoue les commits sur une autre base pour garder un historique linéaire.
    \item \textbf{Cardinalités Merise (0,1 / 1,n) ?} $\rightarrow$ Elles indiquent le nombre minimum et maximum de participations d'une entité à une association ; ex. « un contribuable dépose (0,n) déclarations, une déclaration appartient à (1,1) contribuable ».
    \item \textbf{Algorithmes d'ordonnancement du processeur ?} $\rightarrow$ FIFO (premier arrivé, premier servi), SJF (\textit{Shortest Job First}, le plus court d'abord), Round Robin (tourniquet : chaque processus reçoit un quantum de temps).
    \item \textbf{Qu'est-ce qu'une matrice RACI ?} $\rightarrow$ Tableau des rôles dans un projet : \textbf{R}éalise, \textbf{A}pprouve (responsable final), \textbf{C}onsulté, \textbf{I}nformé.
    \item \textbf{Qu'est-ce qu'un schéma directeur informatique ?} $\rightarrow$ Plan stratégique pluriannuel qui aligne les projets informatiques sur les objectifs de l'administration (priorités, budget, calendrier).
    \item \textbf{PMBOK ?} $\rightarrow$ \textit{Project Management Body of Knowledge} : référentiel international de gestion de projet (PMI).
    \item \textbf{Loi 53-05 ?} $\rightarrow$ Loi de 2007 sur l'échange électronique de données juridiques (valeur de l'écrit électronique), complétée pour la signature électronique par la \textbf{loi 43-20} sur les services de confiance.
    \item \textbf{Qu'est-ce que maCERT ?} $\rightarrow$ Le centre de la DGSSI qui publie les alertes de sécurité et aide les organismes à répondre aux incidents (CERT = \textit{Computer Emergency Response Team}).
\end{enumerate}

\subsection*{I bis. Économie et Statistique (Culture Générale de l'Ingénieur)}
\begin{enumerate}[leftmargin=*]
    \item \textbf{Les trois optiques de calcul du PIB ?} $\rightarrow$ Production (somme des valeurs ajoutées) ; Demande : \textbf{C + I + G + X $-$ M} (C = Consommation des ménages, I = Investissement, G = dépenses publiques (\textit{Government}), X = eXportations, M = iMportations) ; Revenus (salaires + profits + impôts sur la production).
    \item \textbf{PIB nominal ou réel ?} $\rightarrow$ Le PIB réel est corrigé de l'inflation (calculé à prix constants) ; c'est lui qui mesure la croissance.
    \item \textbf{Solde primaire ?} $\rightarrow$ Le solde budgétaire hors intérêts de la dette (il montre si l'État dépense plus qu'il ne perçoit, sans compter le coût de ses emprunts passés).
    \item \textbf{Moyenne ou médiane pour des revenus très inégaux ?} $\rightarrow$ La \textbf{médiane} (valeur qui partage la population en deux moitiés), moins sensible aux très hauts revenus.
    \item \textbf{Indice de Gini ?} $\rightarrow$ Mesure des inégalités : 0 = égalité parfaite, 1 = inégalité maximale.
    \item \textbf{Corrélation proche de 0 ?} $\rightarrow$ Absence de relation \textbf{linéaire}. Piège : une forte corrélation ne prouve pas la causalité.
    \item \textbf{Élasticité-prix de la demande ?} $\rightarrow$ Variation en \% de la quantité demandée quand le prix varie de 1 \%.
    \item \textbf{Effet d'éviction ?} $\rightarrow$ L'emprunt massif de l'État capte l'épargne et renchérit le crédit au détriment de l'investissement privé.
    \item \textbf{Qui produit les statistiques officielles ?} $\rightarrow$ Le \textbf{HCP} (chômage, IPC = Indice des Prix à la Consommation, RGPH = Recensement Général de la Population et de l'Habitat) ; BAM pour la monnaie ; le MEF (TGR, DTFE) pour les finances publiques.
    \item \textbf{Intervalle de confiance à 95 \% ?} $\rightarrow$ Intervalle construit de sorte que, si l'on répétait l'enquête un grand nombre de fois, 95 \% de ces intervalles contiendraient la vraie valeur.
\end{enumerate}

\subsection*{J. Réseaux et Télécommunications Approfondis (Spécialité à 10 postes)}
\begin{enumerate}[leftmargin=*]
    \item \textbf{Quelle est la différence entre un hub, un switch et un routeur ?} $\rightarrow$ Le hub répète le signal vers tous les ports (couche 1, obsolète) ; le switch commute les trames selon l'adresse MAC (couche 2) ; le routeur achemine les paquets entre réseaux selon l'adresse IP (couche 3).
    \item \textbf{Qu'est-ce qu'un switch de niveau 3 ?} $\rightarrow$ Un commutateur capable de router entre VLAN (routage inter-VLAN) sans passer par un routeur externe.
    \item \textbf{À quoi sert le protocole STP ?} $\rightarrow$ \textit{Spanning Tree Protocol} : éviter les boucles de commutation en désactivant les liens redondants, réactivés en cas de panne.
    \item \textbf{Qu'est-ce que le trunking 802.1Q ?} $\rightarrow$ Le transport de plusieurs VLAN sur un même lien entre deux switches, chaque trame portant une étiquette (tag) de VLAN.
    \item \textbf{Que fait le NAT et quelle est sa limite ?} $\rightarrow$ Il traduit des adresses privées en une adresse publique (PAT : plusieurs machines derrière une IP grâce aux ports) ; il complique l'accès entrant et certaines applications (VoIP), d'où le passage progressif à IPv6.
    \item \textbf{Différence entre routage statique et dynamique ; entre OSPF et BGP ?} $\rightarrow$ Statique : routes fixées à la main (petits sites) ; dynamique : échange automatique. OSPF (état de liens, métrique de coût) à l'intérieur d'un système autonome ; BGP (vecteur de chemin) entre systèmes autonomes sur Internet.
    \item \textbf{Qu'est-ce qu'un VPN site à site et un VPN nomade ?} $\rightarrow$ Site à site : tunnel permanent entre deux passerelles (consulat $\leftrightarrow$ ministère), souvent IPsec ; nomade : un utilisateur distant se connecte au réseau via un client VPN (SSL/TLS ou IPsec) avec authentification forte.
    \item \textbf{IPsec : à quoi servent AH, ESP, IKE ?} $\rightarrow$ AH (\textit{Authentication Header}) : intégrité et authentification ; ESP (\textit{Encapsulating Security Payload}) : chiffrement et intégrité ; IKE (\textit{Internet Key Exchange}) : négociation des clés et des associations de sécurité.
    \item \textbf{Qu'est-ce que le MPLS et pourquoi les administrations l'utilisent-elles ?} $\rightarrow$ Une technologie opérateur qui commute par étiquettes et offre des réseaux privés virtuels de niveau 3 (VPN MPLS) avec qualité de service garantie entre sites ; alternative ou complément : SD-WAN sur Internet avec chiffrement.
    \item \textbf{Qu'est-ce que le SD-WAN ?} $\rightarrow$ \textit{Software-Defined WAN} : pilotage logiciel et centralisé des liaisons d'un site (fibre, 4G/5G, MPLS), avec bascule automatique et priorisation des applications ; adapté à des postes à l'étranger aux liaisons hétérogènes.
    \item \textbf{Qu'est-ce que la QoS et quels mécanismes la mettent en œuvre ?} $\rightarrow$ La qualité de service priorise les flux sensibles (voix, visioconférence) : marquage DSCP, files d'attente prioritaires, limitation de débit. Objectifs VoIP : latence $<$ 150 ms, gigue $<$ 30 ms, perte $<$ 1 \%.
    \item \textbf{Comment fonctionne la VoIP ?} $\rightarrow$ La voix est numérisée, compressée par un codec (G.711, G.729), transportée en RTP sur UDP ; SIP établit et termine les appels ; un IPBX gère les postes. Sécurité : SIP sur TLS et SRTP.
    \item \textbf{Différence entre Wi-Fi 5, Wi-Fi 6 et les bandes 2,4 / 5 / 6 GHz ?} $\rightarrow$ Wi-Fi 5 (802.11ac, 5 GHz) ; Wi-Fi 6 (802.11ax, 2,4 et 5 GHz, plus efficace en environnement dense, OFDMA) ; Wi-Fi 6E ajoute la bande 6 GHz. 2,4 GHz porte plus loin mais est plus encombré.
    \item \textbf{Sécurité Wi-Fi : WPA2-Personnel ou WPA2/WPA3-Entreprise ?} $\rightarrow$ Personnel : clé partagée (PSK) ; Entreprise : authentification individuelle via 802.1X et un serveur RADIUS (recommandée en administration), avec un SSID invité isolé.
    \item \textbf{Qu'est-ce que 802.1X ?} $\rightarrow$ Un contrôle d'accès au port réseau (filaire ou Wi-Fi) : le poste doit s'authentifier (certificat ou identifiants) auprès d'un serveur RADIUS avant d'accéder au réseau.
    \item \textbf{Fibre monomode ou multimode ?} $\rightarrow$ Monomode : cœur fin, longues distances (dizaines de km), liaisons opérateur ; multimode : cœur large, courtes distances (bâtiment, salle serveurs). Cuivre : catégories 5e (1 Gbit/s), 6 et 6A (10 Gbit/s sur 100 m).
    \item \textbf{Que signifient FTTH, FTTO, xDSL ?} $\rightarrow$ Fibre jusqu'au domicile / jusqu'au bureau (débit garanti, SLA) ; xDSL : haut débit sur ligne téléphonique en cuivre (ADSL asymétrique, VDSL).
    \item \textbf{Quels sont les principaux apports de la 5G ?} $\rightarrow$ Très haut débit (eMBB), faible latence (URLLC), objets connectés massifs (mMTC), découpage du réseau (\textit{network slicing}). Au Maroc : lancement commercial le 7 novembre 2025, objectif 70 \% de couverture en 2030.
    \item \textbf{Qu'est-ce qu'une liaison VSAT et quand l'utiliser ?} $\rightarrow$ Une liaison satellitaire par petite antenne : secours ou accès principal pour un poste isolé sans infrastructure terrestre fiable ; latence élevée (environ 600 ms en géostationnaire), plus faible avec les constellations en orbite basse.
    \item \textbf{Comment garantir la haute disponibilité d'un service ?} $\rightarrow$ Redondance des équipements (deux pare-feu en cluster, deux liaisons, deux alimentations), répartition de charge, protocoles de passerelle redondante (VRRP/HSRP), sauvegardes, supervision, PRA/PCA avec RTO et RPO définis.
    \item \textbf{Que mesure un SLA de 99,9 \% ?} $\rightarrow$ Une indisponibilité maximale d'environ 8 h 45 par an (99,99 \% : environ 53 minutes).
    \item \textbf{Quels outils de supervision réseau connaissez-vous ?} $\rightarrow$ SNMP (\textit{Simple Network Management Protocol}) pour interroger les équipements, NetFlow pour analyser les flux, syslog pour les journaux ; outils : Zabbix, Nagios, Centreon, PRTG.
    \item \textbf{Qu'est-ce qu'une DMZ et que met-on dedans ?} $\rightarrow$ Une zone tampon entre Internet et le réseau interne, isolée par pare-feu : serveurs web publics, relais de messagerie, reverse proxy ; jamais les bases de données internes.
    \item \textbf{Qu'est-ce qu'un pare-feu nouvelle génération (NGFW) ?} $\rightarrow$ Un pare-feu qui filtre aussi au niveau applicatif : identification des applications et des utilisateurs, IPS intégré, inspection TLS, filtrage d'URL.
    \item \textbf{Qu'est-ce qu'un reverse proxy et un WAF ?} $\rightarrow$ Le reverse proxy reçoit les requêtes des internautes et les transmet aux serveurs internes (masquage, répartition, cache) ; le WAF (\textit{Web Application Firewall}) filtre les attaques web (injections, XSS) devant une application comme rdv.consulat.ma.
    \item \textbf{Comment sécuriser le DNS et la messagerie d'une administration ?} $\rightarrow$ DNSSEC (authenticité des réponses DNS) ; pour la messagerie : SPF, DKIM et DMARC contre l'usurpation d'adresse, chiffrement TLS, filtrage anti-hameçonnage.
    \item \textbf{Qu'est-ce que l'IPv6 et comment migrer ?} $\rightarrow$ Adresses sur 128 bits (notation hexadécimale, ex. 2001:db8::1), plus de NAT nécessaire, auto-configuration (SLAAC) ; migration en double pile (IPv4 et IPv6 simultanés) ou par tunnels.
    \item \textbf{Comment calcule-t-on rapidement un masque ?} $\rightarrow$ /n signifie n bits à 1. Hôtes utilisables $= 2^{(32-n)} - 2$. /25 : 126 ; /26 : 62 ; /27 : 30 ; /28 : 14 ; /29 : 6 ; /30 : 2 (liaison point à point).
    \item \textbf{Qu'est-ce que le VLSM et le CIDR ?} $\rightarrow$ VLSM (\textit{Variable Length Subnet Mask}) : masques de longueur variable pour découper un réseau selon les besoins de chaque sous-réseau ; CIDR (\textit{Classless Inter-Domain Routing}) : notation /n sans classes, agrégation des routes.
    \item \textbf{Différence entre latence, bande passante et débit ?} $\rightarrow$ Latence : délai d'acheminement (ms) ; bande passante : capacité maximale du lien (bit/s) ; débit : quantité réellement transmise.
    \item \textbf{Qu'est-ce qu'un annuaire LDAP / Active Directory et pourquoi le centraliser ?} $\rightarrow$ Une base des utilisateurs, groupes et droits ; centraliser permet l'authentification unique (SSO), l'application des politiques de sécurité et la révocation immédiate des comptes, y compris pour les agents à l'étranger.
    \item \textbf{Qu'est-ce qu'un bastion d'administration ?} $\rightarrow$ Un serveur unique et surveillé par lequel passent obligatoirement toutes les connexions d'administration (SSH, RDP) vers les équipements, avec traçabilité et MFA.
    \item \textbf{Qu'est-ce qu'un PABX/IPBX et un centre d'appel ?} $\rightarrow$ L'autocommutateur téléphonique de l'organisation (IPBX sur IP) ; un centre d'appel s'appuie dessus avec distribution automatique des appels, serveur vocal interactif et supervision, comme le centre d'appel multilingue des services consulaires.
    \item \textbf{Comment relier un poste à l'étranger au ministère de manière sûre et résiliente ?} $\rightarrow$ Deux accès Internet de fournisseurs différents (fibre + 4G/5G ou VSAT en secours), passerelle VPN IPsec redondante vers Rabat, chiffrement de bout en bout pour les flux sensibles, segmentation par VLAN, authentification 802.1X, supervision centralisée, sauvegardes locales et distantes.
\end{enumerate}

\subsection*{K. Exercices Rédigés Corrigés (Contexte du Ministère)}

\begin{tcolorbox}[colback=green!3!white,colframe=green!40!black,title=\textbf{Exercice 1 --- Plan d'adressage d'une ambassade (réseaux)}]
\textbf{Énoncé.} L'ambassade du Maroc dans une capitale européenne dispose du bloc privé \texttt{10.45.0.0/22}. Elle doit héberger : 60 postes administratifs, 25 postes du service consulaire, 12 téléphones IP, 8 serveurs, 20 clients Wi-Fi invités et 2 liaisons point à point vers les équipements de sécurité. Proposez un découpage en sous-réseaux (VLSM), donnez pour chacun l'adresse de réseau, le masque, la plage utilisable et l'adresse de diffusion, puis indiquez les règles de sécurité essentielles entre ces réseaux.

\textbf{Corrigé.}
Le bloc /22 contient $2^{10} = 1024$ adresses (10.45.0.0 à 10.45.3.255). On alloue du plus grand besoin au plus petit :
\begin{itemize}
    \item \textbf{VLAN 10 Administratif} (60 postes $\rightarrow$ 64 adresses, /26) : réseau 10.45.0.0/26, masque 255.255.255.192, plage 10.45.0.1 -- 10.45.0.62, diffusion 10.45.0.63.
    \item \textbf{VLAN 20 Consulaire} (25 postes $\rightarrow$ 32 adresses, /27) : 10.45.0.64/27, masque 255.255.255.224, plage 10.45.0.65 -- 10.45.0.94, diffusion 10.45.0.95.
    \item \textbf{VLAN 30 Wi-Fi invités} (20 $\rightarrow$ 32 adresses, /27) : 10.45.0.96/27, plage 10.45.0.97 -- 10.45.0.126, diffusion 10.45.0.127.
    \item \textbf{VLAN 40 Téléphonie IP} (12 $\rightarrow$ 16 adresses, /28) : 10.45.0.128/28, masque 255.255.255.240, plage 10.45.0.129 -- 10.45.0.142, diffusion 10.45.0.143.
    \item \textbf{VLAN 50 Serveurs} (8 $\rightarrow$ 16 adresses, /28) : 10.45.0.144/28, plage 10.45.0.145 -- 10.45.0.158, diffusion 10.45.0.159.
    \item \textbf{Liaisons point à point} (2 adresses chacune, /30) : 10.45.0.160/30 (plage .161 -- .162, diffusion .163) et 10.45.0.164/30 (plage .165 -- .166, diffusion .167).
\end{itemize}
Le reste du bloc (10.45.0.168 à 10.45.3.255) est réservé à la croissance. Règles de sécurité : le Wi-Fi invités n'accède qu'à Internet (aucun flux vers les autres VLAN) ; le VLAN consulaire accède aux serveurs sur les seuls ports applicatifs nécessaires (HTTPS, base de données via l'application) ; l'administration des équipements passe par un bastion ; la téléphonie est isolée et priorisée par QoS ; tous les flux vers Rabat transitent par le tunnel IPsec ; journalisation centralisée.
\end{tcolorbox}

\begin{tcolorbox}[colback=green!3!white,colframe=green!40!black,title=\textbf{Exercice 2 --- Architecture d'interconnexion d'un consulat avec Rabat (réseaux)}]
\textbf{Énoncé.} Un consulat général reçoit 400 usagers par jour. Il utilise l'application eConsulat hébergée à Rabat, la téléphonie IP et la visioconférence. Proposez une architecture d'interconnexion résiliente et sécurisée, en justifiant chaque composant.

\textbf{Corrigé (plan de réponse attendu).}
\begin{enumerate}
    \item \textbf{Accès :} deux liaisons de fournisseurs différents (fibre professionnelle avec SLA + 4G/5G ou VSAT en secours) ; routeur SD-WAN ou deux routeurs avec bascule automatique.
    \item \textbf{Sécurité périmétrique :} pare-feu nouvelle génération en haute disponibilité (deux boîtiers), IPS activé, filtrage sortant, DMZ locale si un service doit être exposé.
    \item \textbf{Tunnel vers Rabat :} VPN IPsec site à site redondant (deux passerelles côté ministère), chiffrement AES-256, IKEv2, éventuellement VPN MPLS de l'opérateur en complément.
    \item \textbf{Réseau local :} VLAN par usage (guichets, back-office, téléphonie, invités, équipements), 802.1X avec RADIUS, Wi-Fi Entreprise, switches empilés avec liens redondants (STP).
    \item \textbf{Qualité de service :} priorisation de la voix et de la visioconférence (DSCP), réservation de bande passante pour eConsulat.
    \item \textbf{Continuité :} mode dégradé local (cache ou file d'attente si Rabat est injoignable), onduleur, sauvegardes, PRA avec RTO de 4 heures et RPO d'une heure par exemple.
    \item \textbf{Exploitation :} supervision SNMP/NetFlow centralisée à Rabat, journaux vers le SIEM, bastion d'administration, mises à jour planifiées hors heures d'ouverture (attention au décalage horaire).
    \item \textbf{Conformité :} DNSSI, loi 05-20, protection des données personnelles (loi 09-08) ; hébergement des données au Maroc.
\end{enumerate}
\end{tcolorbox}

\begin{tcolorbox}[colback=green!3!white,colframe=green!40!black,title=\textbf{Exercice 3 --- Conception d'un système de rendez-vous consulaire (développement et bases de données)}]
\textbf{Énoncé.} Concevoir la base de données d'une plateforme de prise de rendez-vous consulaire : un usager (identifié par sa CIN ou son passeport) prend un rendez-vous pour un service donné (passeport, état civil, légalisation) dans un consulat, sur un créneau horaire. Un créneau appartient à un consulat et a une capacité. Donnez le MCD (entités, associations, cardinalités), le MLD (tables) et deux requêtes SQL : (a) le nombre de rendez-vous par service pour un consulat donné sur un mois ; (b) les créneaux encore disponibles d'un consulat pour une date.

\textbf{Corrigé.}
\textbf{MCD :} entités USAGER (id, cin\_ou\_passeport, nom, prénom, email, téléphone), CONSULAT (id, ville, pays, fuseau), SERVICE (id, libellé, durée), CRENEAU (id, date, heure\_début, capacité) ; associations : CONSULAT (1,n) --- propose --- (1,1) CRENEAU ; USAGER (0,n) --- prend --- (1,1) RENDEZ\_VOUS (1,1) --- concerne --- (0,n) SERVICE ; RENDEZ\_VOUS (1,1) --- occupe --- (0,n) CRENEAU. Le rendez-vous est une entité (avec statut : confirmé, annulé, honoré) car il porte des attributs propres.

\textbf{MLD :} \texttt{usager(id\_usager PK, piece\_identite UNIQUE, nom, prenom, email, telephone)} ; \texttt{consulat(id\_consulat PK, ville, pays, fuseau)} ; \texttt{service(id\_service PK, libelle, duree\_min)} ; \texttt{creneau(id\_creneau PK, id\_consulat FK, date\_creneau, heure\_debut, capacite)} ; \texttt{rendez\_vous(id\_rdv PK, id\_usager FK, id\_service FK, id\_creneau FK, statut, date\_creation)}. Contrainte : le nombre de rendez-vous confirmés par créneau ne dépasse pas la capacité (vérification applicative ou trigger). Index sur \texttt{(id\_consulat, date\_creneau)}.

\textbf{Requête (a) :}
\begin{verbatim}
SELECT s.libelle, COUNT(*) AS nb_rdv
FROM rendez_vous r
JOIN creneau c ON c.id_creneau = r.id_creneau
JOIN service s ON s.id_service = r.id_service
WHERE c.id_consulat = 12
  AND c.date_creneau BETWEEN '2026-09-01' AND '2026-09-30'
  AND r.statut = 'confirme'
GROUP BY s.libelle
ORDER BY nb_rdv DESC;
\end{verbatim}

\textbf{Requête (b) :}
\begin{verbatim}
SELECT c.id_creneau, c.heure_debut,
       c.capacite - COUNT(r.id_rdv) AS places_restantes
FROM creneau c
LEFT JOIN rendez_vous r
       ON r.id_creneau = c.id_creneau AND r.statut = 'confirme'
WHERE c.id_consulat = 12 AND c.date_creneau = '2026-10-05'
GROUP BY c.id_creneau, c.heure_debut, c.capacite
HAVING c.capacite - COUNT(r.id_rdv) > 0
ORDER BY c.heure_debut;
\end{verbatim}
\textbf{Points de sécurité à mentionner :} requêtes paramétrées (anti-injection), authentification de l'usager (code envoyé par SMS ou e-mail), limitation du nombre de rendez-vous par pièce d'identité (anti-revente de créneaux, CAPTCHA, limitation de débit), journalisation, chiffrement TLS, données personnelles minimisées et conservées selon la loi 09-08.
\end{tcolorbox}

\begin{tcolorbox}[colback=green!3!white,colframe=green!40!black,title=\textbf{Exercice 4 --- Architecture applicative et API (développement)}]
\textbf{Énoncé.} Proposez l'architecture logicielle de la plateforme de rendez-vous (couches, technologies, API) et décrivez le cycle de développement et de déploiement.

\textbf{Corrigé (plan de réponse).} Architecture en trois tiers : interface web responsive (et application mobile) $\rightarrow$ API REST sécurisée (HTTPS, authentification par jetons, limitation de débit) $\rightarrow$ base relationnelle (PostgreSQL) avec sauvegardes. Exemple de ressources REST : \texttt{GET /consulats/\{id\}/creneaux?date=2026-10-05}, \texttt{POST /rendez-vous}, \texttt{DELETE /rendez-vous/\{id\}} ; codes 201 (créé), 409 (créneau complet), 401/403 (accès refusé). Séparation MVC côté application, file de messages pour les notifications (SMS, e-mail), cache pour les créneaux. Cycle : dépôt Git, revues de code, tests unitaires et d'intégration, pipeline CI/CD, environnements de recette et de production, conteneurs (Docker) orchestrés, supervision et journaux centralisés, gestion des incidents selon ITIL. Sécurité : OWASP Top 10, tests d'intrusion avant mise en production, hébergement souverain, conformité DNSSI et loi 09-08.
\end{tcolorbox}

\begin{tcolorbox}[colback=green!3!white,colframe=green!40!black,title=\textbf{Exercice 5 --- Gestion d'un incident de sécurité (cybersécurité)}]
\textbf{Énoncé.} Un agent d'un consulat signale qu'il a cliqué sur un lien reçu par courriel et que ses fichiers deviennent illisibles. Décrivez les actions à mener dans l'heure, dans la journée, puis les mesures durables.

\textbf{Corrigé.} \textbf{Dans l'heure :} isoler le poste du réseau (débrancher, désactiver le Wi-Fi) sans l'éteindre pour préserver les preuves ; alerter la DSI et le responsable sécurité ; identifier les partages réseau accessibles depuis le poste et les mettre hors ligne ; changer les mots de passe de l'agent. \textbf{Dans la journée :} qualifier l'incident (rançongiciel), rechercher d'autres postes touchés (indicateurs de compromission), bloquer le domaine et les adresses malveillantes sur le pare-feu et la messagerie, informer maCERT/DGSSI conformément à la loi 05-20, restaurer depuis des sauvegardes saines après vérification, communiquer en interne, ne pas payer de rançon. \textbf{Durablement :} sauvegardes 3-2-1 testées, mises à jour, EDR sur les postes, MFA, segmentation, filtrage des pièces jointes, sensibilisation régulière des agents, exercice de gestion de crise, revue des droits d'accès (moindre privilège), notification à la CNDP si des données personnelles sont concernées.
\end{tcolorbox}

\begin{tcolorbox}[colback=green!3!white,colframe=green!40!black,title=\textbf{Exercice 6 --- Algorithmique (développement)}]
\textbf{Énoncé.} Écrire en pseudo-code une fonction qui, à partir d'une liste de rendez-vous (heure de début, durée en minutes) triés par heure, détecte les chevauchements et retourne leur nombre.

\textbf{Corrigé.}
\begin{verbatim}
FONCTION compterChevauchements(rdv : liste triée par debut) : entier
    nb <- 0
    finPrecedente <- -infini
    POUR i DE 1 A taille(rdv) FAIRE
        debut <- rdv[i].debut
        fin   <- rdv[i].debut + rdv[i].duree
        SI debut < finPrecedente ALORS
            nb <- nb + 1
        FIN SI
        finPrecedente <- max(finPrecedente, fin)
    FIN POUR
    RETOURNER nb
FIN FONCTION
\end{verbatim}
Complexité : $O(n)$ après tri (le tri lui-même coûte $O(n \log n)$). Piège classique : comparer chaque rendez-vous à tous les autres donne $O(n^2)$.
\end{tcolorbox}
